\chapter{Reference: \code{run} and \code{RunConfig}}
\label{ch:refman-run}

\begin{aside}
\code{run<P>} is the front door. It hands control of stdin,
stdout, and the terminal modes to \maya{}, executes your
program, then politely puts everything back.
\end{aside}

\section{Signature}

\begin{lstlisting}[language=C++]
template <Program P>
int run(RunConfig cfg = {});
\end{lstlisting}

\noindent Returns the process exit code (\code{0} on
graceful quit).

\section{RunConfig fields}

\begin{itemize}[leftmargin=*]
  \item \code{title}: window title, where supported.
  \item \code{alt\_screen}: switch to the alternate screen
    on start, restore on exit. Default \code{true}.
  \item \code{mouse}: enable mouse reporting. Default
    \code{false}.
  \item \code{fps}: target frame rate. Default \code{60}.
  \item \code{cursor}: \code{Hidden}, \code{Block}, or
    \code{Bar}. Default \code{Hidden}.
\end{itemize}

\section{Working example}

\begin{lstlisting}[language=C++]
#include <maya/maya.hpp>

using namespace maya;
using namespace maya::dsl;

struct Demo {
    struct Model { int frames = 0; };
    struct Tick {};
    struct Quit {};
    using Msg = std::variant<Tick, Quit>;

    static Model init() { return {}; }

    static auto update(Model m, Msg msg)
        -> std::pair<Model, Cmd<Msg>>
    {
        return std::visit(overload{
            [&](Tick) { m.frames++; return std::pair{m, Cmd<Msg>{}}; },
            [](Quit)  { return std::pair{Model{}, Cmd<Msg>::quit()}; },
        }, msg);
    }

    static Element view(const Model& m) {
        return v(
            t<"run() demo"> | Bold,
            blank_,
            text(std::format("frames: {}", m.frames)),
            blank_,
            t<"q to quit"> | Dim
        ) | pad<1> | border_<Round>;
    }

    static auto subscribe(const Model&) -> Sub<Msg> {
        return Sub<Msg>::batch(
            key_map<Msg>({ {'q', Quit{}} }),
            Sub<Msg>::every(std::chrono::milliseconds(16), Tick{})
        );
    }
};

int main() {
    return run<Demo>({
        .title      = "demo",
        .alt_screen = true,
        .mouse      = false,
        .fps        = 60,
    });
}
\end{lstlisting}

\section{Notes}

\begin{itemize}[leftmargin=*]
  \item Returning from \code{main} after \code{run} is fine;
    the runtime restores the terminal on its own.
  \item If \code{alt\_screen} is \code{false}, output stays
    in scrollback. Good for tools whose output you want to
    \code{grep} later.
  \item \code{fps} caps the loop, it does not force frames.
    Idle frames are skipped.
  \item Setting \code{mouse = true} is necessary before any
    mouse event will reach your \code{subscribe}.
\end{itemize}

\section{Recap}

\begin{recap}
\begin{itemize}[leftmargin=*]
  \item One function, one config struct.
  \item Title, alt-screen, mouse, fps, cursor.
  \item Returns an exit code.
\end{itemize}
\end{recap}

\section*{Workshop}

\begin{enumerate}[leftmargin=*]
  \item Flip \code{alt\_screen} to \code{false}. Watch the
    UI render in scrollback.
  \item Lower \code{fps} to \code{5}. Notice the change.
  \item Enable \code{mouse} and add a mouse handler in
    \code{subscribe}.
\end{enumerate}
